\section{Limitations}
\label{sec:limitations}

The study has several important limits. First, the frozen benchmark contains only 176 accepted trajectories from one GQA-derived subset and one primary VLM family/model. Questions are concentrated in color, identity, material, and shape categories. Evidence regions are scene-graph object boxes rather than exhaustive human rationales, and the main audit used automated visual quality control with only a limited human spot-check. Strict normalized exact matching treats semantically close responses such as ``tan'' and ``khaki'' as different; we did not add an unreported semantic-equivalence analysis.

Only local Gaussian blur is tested strongly as an unseen degradation family. The method does not improve AURC over native confidence, and learned-head AUROC/Brier differences are not statistically significant. Order constraints encode an epistemic prior that may fail for individual samples, particularly when correlated context remains. Matched non-critical boxes are equally sized and geometrically separated but can overlap other annotated objects; only 32/170 have zero annotated-object overlap. Category labels were corrected after initial processing in P17, so stale category metadata in some blur caches must not be used for category-specific claims. Finally, a larger study should test multiple VLM families, datasets, evidence annotations, and intervention mechanisms before generalizing beyond this controlled setting.
